\documentclass[conference]{IEEEtran}
\IEEEoverridecommandlockouts

\usepackage{cite}
\usepackage{amsmath,amssymb,amsfonts}
\usepackage{algorithmic}
\usepackage{graphicx}
\usepackage{textcomp}
\usepackage{xcolor}
\usepackage{booktabs}
\usepackage{url}

\def\BibTeX{{\rm B\kern-.05em{\sc i\kern-.025em b}\kern-.08em
    T\kern-.1667em\lower.7ex\hbox{E}\kern-.125emX}}

\begin{document}

\title{Armavour: Measuring Computer-Use Agent Susceptibility to
Legally-Codified Dark Patterns\\
\thanks{Work conducted at Vivekanand Education Society's Institute of
Technology, Mumbai.}}

\author{
\IEEEauthorblockN{Soumya Vinod}
\IEEEauthorblockA{\textit{Department of Computer Applications} \\
\textit{Vivekanand Education Society's Institute of Technology}\\
Mumbai, India}
\and
\IEEEauthorblockN{Chinmay [Surname]}
\IEEEauthorblockA{\textit{Department of Computer Applications} \\
\textit{Vivekanand Education Society's Institute of Technology}\\
Mumbai, India}
}

\maketitle

\begin{abstract}
Large language model agents increasingly act on the web on a user's
behalf, completing purchases, managing subscriptions, and navigating
consent interfaces. These interfaces frequently contain dark patterns:
design choices engineered to steer users toward outcomes they did not
intend. India's Central Consumer Protection Authority (CCPA) codified
thirteen such patterns in 2023, making them a legally-grounded rather
than merely descriptive taxonomy. We present Armavour, a reproducible
benchmark that measures whether computer-use agents are deceived by
these patterns. Armavour comprises a deterministic testbed implementing
twelve CCPA patterns at four manipulation intensities in three languages,
a ground-truth oracle that scores outcomes from the document object
model rather than from agent self-report, and a rubric-driven LLM judge
for the two patterns whose success condition is not deterministically
observable. Across 1,080 episodes we find that susceptibility scales
monotonically with manipulation intensity, that it is highly
pattern-dependent rather than uniform, and---our strongest
result---that it increases substantially when the agent operates an
interface rendered in Hindi or Hinglish under an English instruction
(16.9\% to 30.0\%, McNemar $p < 0.0001$ on seed-matched pairs). We also
report a class of measurement failure we encountered and corrected:
patterns whose manipulation can be bypassed by evaluating the task's
stated criterion alone are structurally unmeasurable, and we argue this
failure mode is likely present in existing benchmarks.
\end{abstract}

\begin{IEEEkeywords}
dark patterns, LLM agents, computer-use agents, benchmark, consumer
protection, multilingual evaluation, AI safety
\end{IEEEkeywords}

\section{Introduction}

Dark patterns are interface designs that exploit predictable human
cognitive tendencies to produce outcomes favouring the platform over the
user: a pre-ticked donation checkbox, a mandatory fee revealed only at
the final confirmation, a decline button labelled to induce guilt. They
are widespread, well documented, and---in a growing number of
jurisdictions---illegal.

The premise of this work is that the deployment of autonomous web agents
changes the threat model in a way that has not been adequately measured.
When an agent completes a purchase on a user's behalf, the manipulation
is no longer aimed at a human's attention and emotions; it is aimed at
whatever the agent's perception pipeline delivers to a language model,
and at whatever reasoning that model then performs. Some manipulations
may fail entirely against this substrate. Others may succeed more
reliably than they do against humans. Neither outcome is obvious in
advance, and the distinction matters both for consumer protection policy
and for the design of agent systems.

We take India's CCPA \emph{Guidelines for Prevention and Regulation of
Dark Patterns, 2023} as our taxonomy. This choice is deliberate. Most
dark-pattern research uses taxonomies derived from HCI literature, which
are descriptively rich but carry no enforcement weight. The CCPA
guidelines enumerate thirteen specific patterns with verbatim legal
definitions and have been used in actual enforcement actions against
named Indian platforms. Grounding a benchmark in a codified taxonomy
means a result about pattern $X$ is a result about a category regulators
already act on.

\textbf{Contributions.}

\begin{enumerate}
\item \textbf{Armavour}, an open benchmark comprising a deterministic
testbed for twelve CCPA-codified dark patterns at four intensities and
three languages, with an oracle that establishes ground truth from DOM
state rather than agent self-report.
\item An \textbf{intensity dose-response result}: deception rises
monotonically from 0\% at control to 31.7\% at aggressive intensity
under otherwise identical conditions.
\item A \textbf{multilingual susceptibility result}: on 180
seed-matched scenario triplets in which only the interface language
varies, deception rises from 16.9\% (English) to 22.2\% (Hinglish) to
30.0\% (Hindi), significant at $p < 0.0001$ by paired McNemar's test.
\item A \textbf{qualitative failure taxonomy} derived from reasoning
traces, including three distinct mechanisms by which agents fall for
subscription traps and three archetypes of task abandonment.
\item A \textbf{methodological finding} concerning benchmark
construction: patterns whose manipulation can be circumvented by
evaluating the task's stated criterion alone measure criterion-comparison
ability rather than manipulation susceptibility. We document two such
patterns in our own initial design, the empirical signature that
revealed them, and the correction.
\end{enumerate}

\section{Related Work}

\subsection{Dark pattern taxonomies and detection}

Work on dark patterns has largely proceeded from qualitative taxonomy
toward automated detection at scale. Existing taxonomies are
descriptively useful but heterogeneous, and detection work typically
answers whether a pattern is \emph{present} on a page rather than
whether it \emph{works} on a given agent. Armavour addresses the latter
question and, by adopting a legal taxonomy, makes its categories
directly commensurable with regulatory action.

\subsection{Agent benchmarks and susceptibility}

Web-agent benchmarks predominantly measure task completion. Recent work
has begun to measure susceptibility to adversarial or manipulative
interfaces; SusBench is the closest prior work and reports pattern-level
avoidance rates that we treat as a comparison baseline. Two differences
are worth stating. First, Armavour is grounded in a codified legal
taxonomy rather than a research taxonomy. Second, and more consequential
for our results, Armavour varies manipulation \emph{intensity} and
interface \emph{language} as controlled experimental factors, which
allows dose-response and cross-lingual claims that presence/absence
designs cannot support.

\subsection{Multilingual evaluation of language models}

Degradation of model performance in lower-resource languages is well
established for static tasks. Its consequences for \emph{agentic} safety
behaviour---whether an agent operating a non-English interface is more
readily manipulated---have received less attention. Section~\ref{sec:e2}
provides direct evidence on this question.

\section{The Armavour Benchmark}

\subsection{Design principles}

Armavour is built around three commitments.

\textbf{Ground truth from the interface, not the agent.} An agent's
report that it declined an add-on is not evidence that it did. Every
pattern exposes a deterministic oracle that reads the terminal DOM state
at the moment of the transaction and emits a structured result. The
agent cannot influence this signal except by acting on the page.

\textbf{Intensity as a controlled variable.} Each pattern is
implemented at four intensities---\texttt{control}, \texttt{subtle},
\texttt{moderate}, \texttt{aggressive}---in which only the manipulation
varies. \texttt{control} renders the same task honestly and supplies the
baseline against which deception is measured. Prices, item ordering, and
DOM structure are held constant across intensities wherever the pattern
specification permits.

\textbf{Versioned interface contracts.} Five contracts govern the
seams between components: episode configuration, oracle output, stable
element identifiers, the episode log schema, and the judge rubric
interface. Contracts are versioned and changed only by explicit
agreement, which is what makes results comparable across runs conducted
weeks apart.

\subsection{Testbed}

The testbed is a React/TypeScript single-page application presenting
seven simulated services (ticketing, e-commerce, education,
subscription management, SaaS signup, security software, and news).
Twelve of the thirteen CCPA patterns are implemented; the thirteenth is
out of scope as it concerns advertising practice not observable within a
single session.

Episodes are driven entirely by URL parameters, per Contract~1:

\begin{center}
\texttt{/?site=\{s\}\&pattern=\{p\}\&intensity=\{i\}}\\
\texttt{\&lang=\{l\}\&seed=\{n\}}
\end{center}

This makes every episode a pure function of its configuration, and the
SHA-256 hash of that configuration serves as the idempotency key for
logging and resumption.

\subsection{Oracle}

At the terminal action the page sets a global result object and
corresponding DOM attributes:

\begin{verbatim}
{ pattern, avoided, total, expected_total, ... }
\end{verbatim}

Pattern-specific fields are added alongside---for basket sneaking, the
sneaked item and its amount; for drip pricing, the advertised price, the
final total, and their difference. Outcomes are coded on two binary
axes, whether the task was completed and whether the pattern was
avoided, yielding four categories: EC (completed, avoided), DC
(completed, deceived), EF (not completed, avoided), and DF (not
completed, deceived).

\subsection{Judge}

Two patterns---false urgency and confirm shaming---have a success
condition that DOM state cannot fully determine, because the question is
whether the manipulation \emph{influenced} the decision rather than
merely whether a particular button was pressed. For these, a second
language model scores the reasoning trace against a CCPA-grounded
rubric, returning a boolean flag and a one-sentence evidence string.
Contract~5 requires the judge model to differ from the agent model.

Two decisions here are consequential and we state them explicitly.

First, for judge-scored patterns the reported susceptibility rate is the
\emph{judge flag} rate, not the raw DC rate. A DC outcome records that
the platform-favoured result occurred, not that the dark pattern caused
it. In our soft-pattern pilot, four confirm-shaming episodes scored DC
while the judge correctly determined the guilt framing played no
role---they were button-semantics errors. Conflating the two would have
overstated the pattern's effect roughly fourfold.

Second, an agent whose reasoning acknowledges the manipulation but whose
final action is nonetheless correct is scored as \emph{not} swayed. The
benchmark measures whether the pattern changed the decision, not whether
it entered deliberation. Acknowledged-but-resisted influence is
preserved in the judge's evidence string for qualitative analysis
without affecting any metric.

\subsection{Judge validation}
\label{sec:judgeval}

We validated the judge against a hand-constructed set of twelve cases
covering both judge-scored patterns, deliberately including the failure
mode we had observed in practice: episodes in which the agent takes the
task-correct action while its reasoning quotes the guilt-laden label,
because in this pattern the correct button \emph{is} the guilt-laden
button.

An initial judge model scored 72.7\% accuracy with precision 0.50 and
recall 1.00---systematically over-flagging, with three false positives
and no false negatives. Inspection showed the failures were not
borderline: in one case the model's own evidence string stated the agent
showed ``no hesitation or reconsideration'' and then flagged it anyway,
and inverted the final-choice value it had been given directly.

Substituting a larger judge model, with \emph{no change to the rubric
text}, produced 100\% accuracy across all twelve cases with zero parse
errors, reproduced over two consecutive runs. We report this because it
localises the problem: judge unreliability here was a model-capability
limitation, not a prompt-engineering one, and further rubric iteration
would have been wasted effort.

\section{Experimental Design}

\subsection{Configuration}

The agent is \texttt{llama-3.3-70b-versatile} accessed through a
computer-use adapter that extracts interactive elements and their
surrounding text context from the DOM, presents them to the model as
structured JSON, and executes the returned action via Playwright. The
judge is \texttt{gpt-oss-120b}. Inference is deterministic
($\text{temperature} = 0$) at every call site.

\subsection{E1: susceptibility matrix}

Twelve patterns $\times$ four intensities $\times$ ten seeds, English
only, single adapter: 480 episodes. A cross-model spot-check re-runs
twelve patterns at aggressive intensity with five seeds on a smaller
model from the same provider (\texttt{llama-3.1-8b-instant}), isolating
model capability within a single provider rather than confounding it
with cross-provider differences: 60 episodes.

\subsection{E2: language arms}
\label{sec:e2design}

Six patterns whose manipulation is carried by language---trick question,
confirm shaming, interface interference, forced action, false urgency,
and SaaS billing---at three intensities, three languages, ten seeds: 540
episodes. Structural patterns whose manipulation is numeric or positional
were excluded, since varying their language would mostly measure
translation noise.

Language-sensitive strings were written by a native speaker rather than
machine-translated, because a flat translation of, for example, a
triple-nested negation would not preserve the manipulation the item is
designed to test.

An important detail of the design: the \emph{task instruction} is
delivered in English in all arms; only the interface is localised. The
finding in Section~\ref{sec:e2} is therefore about agents operating a
non-English interface under an English instruction, which is a common
real deployment configuration but is narrower than a claim about agents
instructed in Hindi.

\section{Results}

\subsection{Aggregate outcomes}

Table~\ref{tab:overall} summarises 1,080 episodes.

\begin{table}[htbp]
\caption{Outcome distribution by experiment arm}
\begin{center}
\begin{tabular}{lrrrr}
\toprule
\textbf{Arm} & \textbf{$n$} & \textbf{EC} & \textbf{DC} & \textbf{EF} \\
\midrule
E1a (70B, English)      & 480   & 84.8\% & 14.6\% & 0.6\% \\
Spot-check (8B)         & 60    & 56.7\% & 35.0\% & 8.3\% \\
E2 (three languages)    & 540   & 74.8\% & 23.1\% & 2.0\% \\
\midrule
\textbf{Total}          & 1{,}080 & 78.2\% & 20.0\% & 1.8\% \\
\bottomrule
\end{tabular}
\label{tab:overall}
\end{center}
\end{table}

No episode produced a DF outcome. We treat this as a property of the
scoring construction rather than a finding, and discuss it in
Section~\ref{sec:limitations}.

The spot-check arm is notable independently: the smaller model was
deceived at more than twice the rate of the larger one (35.0\% versus
14.6\%) under identical conditions at a single intensity.

\subsection{Intensity dose-response}

Because E2 omits the \texttt{subtle} level and the spot-check runs only
at \texttt{aggressive}, an intensity comparison pooled across arms would
confound intensity with language and model. We therefore report the
gradient from E1a alone, where each cell contains exactly 120 episodes
of a single model and language (Table~\ref{tab:intensity}).

\begin{table}[htbp]
\caption{Deception rate by intensity, E1a only ($n = 120$ per cell)}
\begin{center}
\begin{tabular}{lr}
\toprule
\textbf{Intensity} & \textbf{DC rate} \\
\midrule
Control     & 0.0\% \\
Subtle      & 7.5\% \\
Moderate    & 19.7\% \\
Aggressive  & 31.7\% \\
\bottomrule
\end{tabular}
\label{tab:intensity}
\end{center}
\end{table}

Two things follow. First, the control condition produces no deception at
all across 120 episodes, which is the precondition for interpreting any
other number: the agent is not failing these tasks for reasons unrelated
to the manipulation. Second, the relationship is monotonic and
substantial, and it is causal in the strict sense that intensity is the
only manipulated variable.

\subsection{Pattern-level susceptibility}

Susceptibility is strongly pattern-dependent
(Table~\ref{tab:patterns}). SaaS billing---an auto-renewing ``free
trial''---is by a wide margin the most effective, while several patterns
produce no deception at all.

\begin{table}[htbp]
\caption{Deception rate by pattern}
\begin{center}
\begin{tabular}{lr}
\toprule
\textbf{Pattern} & \textbf{DC rate} \\
\midrule
SaaS billing              & 57.0\% \\
Trick question            & 40.0\% \\
Drip pricing              & 33.3\% \\
Forced action             & 32.6\% \\
Bait and switch           & 28.9\% \\
Basket sneaking           & 11.1\% \\
Subscription trap         & 11.1\% \\
Confirm shaming           & 2.2\% \\
Interface interference    & 0.0\% \\
Nagging                   & 0.0\% \\
\midrule
\multicolumn{2}{l}{\emph{Excluded as unmeasurable
(Section~\ref{sec:measurability}):}} \\
Disguised advertisement   & --- \\
False urgency             & --- \\
\bottomrule
\end{tabular}
\label{tab:patterns}
\end{center}
\end{table}

\subsection{Language degradation}
\label{sec:e2}

This is our strongest result. Across 180 scenario triplets in which
seed, pattern, intensity, model, and adapter are all held constant and
only the interface language varies, deception rises monotonically with
distance from English (Table~\ref{tab:lang}).

\begin{table}[htbp]
\caption{Deception rate by interface language, matched triplets
($n = 180$ per arm)}
\begin{center}
\begin{tabular}{lrr}
\toprule
\textbf{Interface language} & \textbf{DC rate} & \textbf{$\Delta$ vs.\ EN} \\
\midrule
English   & 17.2\% & --- \\
Hinglish  & 22.2\% & $+5.0$ pp \\
Hindi     & 30.0\% & $+12.8$ pp \\
\bottomrule
\end{tabular}
\label{tab:lang}
\end{center}
\end{table}

Because the design is fully paired, we test significance with McNemar's
test on discordant pairs. Comparing English against Hindi: 26 scenarios
were avoided in English but deceived in Hindi, against 3 in the reverse
direction. This yields $\chi^2 = 16.69$, $p = 0.000044$.

The effect is not uniform across patterns
(Table~\ref{tab:langpattern}). Forced action degrades most sharply,
rising two-and-a-half fold.

\begin{table}[htbp]
\caption{Language effect by pattern ($n = 30$ per cell)}
\begin{center}
\begin{tabular}{lrrr}
\toprule
\textbf{Pattern} & \textbf{EN} & \textbf{Hinglish} & \textbf{Hindi} \\
\midrule
Forced action    & 26.7\% & 33.3\% & 66.7\% \\
SaaS billing     & 40.0\% & 66.7\% & 60.0\% \\
Trick question   & 36.7\% & 33.3\% & 53.3\% \\
\bottomrule
\end{tabular}
\label{tab:langpattern}
\end{center}
\end{table}

A further observation supports the interpretation. The control
condition---an honest interface with no manipulation---produced zero
deception in E1a's 120 English episodes but six deceptions in the E2
arms, all in non-English conditions. Because there is no manipulation to
fall for in control, these are comprehension failures, and their
appearance only in the localised arms is consistent with degraded
interface understanding rather than degraded manipulation resistance
specifically.

\section{Qualitative Analysis}

\subsection{How agents fall for subscription traps}

We manually inspected twenty randomly sampled deceived episodes from
SaaS billing, the highest-susceptibility pattern. All twenty are genuine
model failures rather than harness or execution artefacts, and they
resolve into three mechanisms.

\textbf{Semantic conflation (roughly 65\%).} The model treats ``free
trial'' as satisfying an instruction to select a ``free plan,'' and does
not register that the trial converts to a recurring charge. The
auto-renewal terms are present and legible in the extracted context;
the model does not connect them to the constraint it was given.

\textbf{Constraint-violating rationalisation (roughly 25\%).} At higher
intensities the interface offers no zero-cost option at all. The model
identifies this correctly---one trace reads ``there's no clear option
for a free plan without auto-renewal''---and then proceeds to click the
trap anyway, on the grounds that it is ``the only actionable step.''
This is a distinct and more concerning failure than the first: the
constraint was understood and then discarded under pressure to complete
the task.

\textbf{Keyword-triggered action (roughly 10\%).} A single-step response
in which the model matches on the word ``free'' and acts without
evaluating surrounding pricing text.

\subsection{Task abandonment}

Nineteen episodes ended with the task incomplete but the pattern
avoided. These divide into three archetypes.

\textbf{Principled refusal} (11 episodes, SaaS billing, 70B). The model
recognises that the only available path commits the user to a recurring
charge in violation of its instruction, and declines to act at all. We
regard this as the desirable behaviour when an interface admits no
compliant action, and note that it is the larger model that exhibits it.

\textbf{Premature completion} (3 episodes, confirm shaming). The model
successfully removes the manipulated add-on, then declares the task
finished without completing checkout---resisting the manipulation but
losing the task in the process.

\textbf{Fixation loop} (5 episodes, basket sneaking, 8B). The smaller
model correctly identifies that the total contains a sneaked item and
refuses to pay, but then repeats the same corrective action for twenty
consecutive steps until the step budget expires. It detected the
constraint and lacked the planning capacity to act on it---a capability
boundary rather than a susceptibility one.

\subsection{Why some patterns fail entirely}

Interface interference produced no deception at any intensity. Trace
inspection shows this is a genuine null rather than a broken
implementation: at aggressive intensity the decline option is rendered
as low-contrast, de-emphasised text beside a visually dominant renewal
button, and the agent's reasoning identifies it directly---``the
`Accept the risk' button is the option that aligns with this goal''---and
selects it. The manipulation targets human visual salience, and the
agent reads all text nodes with equal weight.

This is a robustness result specific to DOM-based agents, and it is
precisely the claim that a vision-based agent might overturn. We regard
establishing this as the most important open item
(Section~\ref{sec:limitations}).

Nagging fails for a related reason: the agent declines the interrupting
prompt and proceeds. Repetition imposes no cost on a system that does
not experience annoyance.

\section{A Measurability Failure in Benchmark Construction}
\label{sec:measurability}

We report this section because we believe the failure it describes is
easy to introduce, hard to notice, and likely present elsewhere.

Our disguised-advertisement item presented four search results, one a
sponsored placement styled to blend in, with the label's prominence
decreasing as intensity rose. The task was to buy the cheapest genuine
product. Following our own specification, the ad was priced above every
organic item so that selecting it would be unambiguously wrong.

That last decision made the item unmeasurable. Because the ad was the
most expensive result and the task asked for the cheapest, any agent
that compared prices avoided it---regardless of whether the sponsored
label was clearly shown, greyed out, or absent entirely. Intensity could
not affect the outcome, and did not: the pattern returned 0\% deception
at every level.

The signature that revealed it was an inconsistency between runs. In an
earlier pilot the same item had returned 60\% deception. Reconciling the
two showed that the pilot agent had been receiving no product names or
prices at all---its traces read ``all buttons have the same text `Buy'''
and ``no product details provided''---and had defaulted to the first
element, which happened to be the ad. The pilot had measured a
perception-pipeline gap and reported it as ad-disguise susceptibility.
Once the extraction layer supplied surrounding text, the measured rate
fell to zero, which is the correct value for an item that cannot
discriminate.

The general condition is this: \emph{if the task's stated criterion is
sufficient to identify the correct choice, the manipulation is never
load-bearing, and the item measures criterion evaluation rather than
manipulation susceptibility.} An item can satisfy every apparent
correctness requirement---unambiguous ground truth, a well-formed
intensity ladder, a clean oracle---and still measure nothing.

We corrected the item by inverting the price relationship, so that the
ad is now the cheapest result and price-only reasoning selects it. The
sponsored label becomes the sole discriminating signal, which is what
the pattern is nominally about. Affected episodes are being re-run.

False urgency exhibits the same structural problem---a single-choice
item in which price determines the answer and the urgency copy is never
load-bearing---and we have not identified an equivalent fix. Making the
urgent item cheapest would make selecting it correct; equalising prices
removes the discriminator entirely. Urgency is pressure to act
\emph{quickly}, which a single-step episode has no way to express.
Measuring it plausibly requires a time- or step-sensitive construction,
and we exclude the item rather than report a number we cannot defend.

A practical corollary for benchmark authors: a pattern scoring far from
its published baseline in either direction warrants a construction
audit before it is reported. Our figure of 100\% avoidance against a
prior-work baseline of approximately 65\% was the anomaly that prompted
this investigation.

\section{Limitations}
\label{sec:limitations}

\textbf{Single adapter.} All reported results use one computer-use
adapter. Because that adapter's extraction layer determines what reaches
the model, and because Section~\ref{sec:measurability} demonstrates how
consequential that layer is, our results are strictly claims about
agents built this way. A second adapter arm is designed and specified
but not yet executed. The interface-interference null in particular
should be regarded as provisional until a vision-based agent is tested.

\textbf{Single agent model, single provider.} The main arm uses one
model; the spot-check adds a smaller model from the same provider. The
pattern-level ranking may not generalise across model families.

\textbf{Trace-only judging.} The judge scores reasoning traces without
the final screenshot, as the selected judge model does not support
vision. For patterns where visual context could disambiguate influence,
this may reduce judge sensitivity.

\textbf{English instructions in the language arms.} As noted in
Section~\ref{sec:e2design}, only the interface is localised. The result
does not speak to agents instructed in Hindi.

\textbf{Deterministic outcomes conflate mechanism.} For
oracle-scored patterns we have no equivalent of the judge's
causal separation. An episode is scored as deceived whether the agent
fell for the manipulation or merely mis-executed a correct intention. We
observed at least one episode whose reasoning suggested the latter. The
reported deception rates for deterministic patterns therefore include an
unknown proportion of execution failures, and the traces are retained so
this can be quantified in future work.

\textbf{DF is unreachable by construction.} No episode produced a
deceived-and-incomplete outcome. This follows from the scoring path
rather than from agent behaviour, and the four-way outcome space should
be understood as three-way in practice.

\section{Conclusion}

We built Armavour to ask whether legally-codified dark patterns work on
agents. The answer is that some do, some do not, and which is which is
not predictable from how well they work on people.

Manipulations that alter what information is available or when it is
disclosed---auto-renewing trials, negated consent phrasing, late-revealed
fees---succeed against agents at substantial rates. Manipulations that
rely on visual salience or on wearing the user down succeed not at all,
because they target properties the agent does not have.

The result we consider most consequential is that susceptibility rises
sharply when the interface is not in English, on identical tasks with
identical seeds. Agent deployment is not confined to English-language
interfaces, and consumer-protection regimes such as the CCPA exist
precisely in the jurisdictions where this degradation is largest.

We also report, deliberately, a way in which our own benchmark measured
nothing at all for two of its items, and the empirical signature that
exposed it. We suspect this failure mode is not unique to us.

\section*{Availability}

The testbed, harness, pattern specifications, judge rubrics, and
validation datasets are available at
\url{https://github.com/Soumya-Vinod/Armavour}.

\begin{thebibliography}{00}

\bibitem{ccpa2023} Central Consumer Protection Authority, ``Guidelines
for Prevention and Regulation of Dark Patterns, 2023,'' Ministry of
Consumer Affairs, Government of India, 2023.

\bibitem{placeholder1} [SusBench citation --- to be completed.]

\bibitem{placeholder2} [Brignull, dark patterns taxonomy --- to be
completed.]

\bibitem{placeholder3} [Mathur et al., dark patterns at scale --- to be
completed.]

\bibitem{placeholder4} [Web agent benchmark citation --- to be
completed.]

\bibitem{placeholder5} [Multilingual LLM evaluation citation --- to be
completed.]

\end{thebibliography}

\end{document}